\section{Experimental Method}
\label{sec:experiments}
All results are available at ~\cite{campArtifact}.

\subsection{Workflows, Datasets, and Platform}

The matched offline cohort contains 33,516 task records from six workflows
(Table~\ref{tab:cohort}). Nextflow/nf-core
\cite{ditommaso2017nextflow,nfcoreFramework} runs on AWS ParallelCluster with
Slurm and 13 worker nodes, each with 32 vCPUs and 256~GiB RAM. %Bowtie2
% \cite{langmead2012bowtie2} is held out from the offline cohort and used only
% for the live cold-to-warm feedback evaluation.

\begin{table}[t]
\caption{Matched workflow cohort.}
\label{tab:cohort}
\centering
\small
\setlength{\tabcolsep}{1.4pt}
\renewcommand{\arraystretch}{1.05}
\begin{tabular}{@{}l l p{0.22\columnwidth} r r r@{}}
\toprule
Workflow & Version & Dataset & All & Test & Proc. \\
\midrule
GATK & 4.5.0.0 custom & Synthetic hg38 reads & 500 & 75 & 1 \\
Minimap2 & custom & ONT-style reads & 6,001 & 900 & 5 \\
Sarek~\cite{nfcoreSarekGithub} & 3.6.0 & GIAB NA12878 WGS & 6,012 & 902 & 25 \\
Seqinsp.~\cite{nfcoreSeqinspectorGithub} & 1.1.0 & ENA metagenomes & 6,001 & 900 & 3 \\
Taxprof.~\cite{nfcoreTaxprofilerGithub} & 2.0.1 & ENA paired-end WGS & 5,999 & 900 & 2 \\
Viralmeta.~\cite{nfcoreViralmetagenomeGithub} & 1.1.3 & ENA metagenomes & 9,003 & 1,349 & 6 \\
\midrule
Total & -- & -- & 33,516 & 5,026 & -- \\
\bottomrule
\end{tabular}
\end{table}

GATK and Minimap2 are deterministic custom workloads constructed for this
study. GATK generates paired reads over hg38 for repeated HaplotypeCaller
scatters, while Minimap2 generates ONT-style reads over fixed genomic windows
and runs mapping, indexing, and summary steps. Sarek uses GIAB NA12878 30x WGS
with GRCh38 interval scattering. Taxprofiler uses 5,998 paired-end Illumina
metagenomic WGS runs from an ENA snapshot. Seqinspector and Viralmetagenome use
3,000 runs from the same ENA cohort with the Kraken2 Standard-8 database.
Together, these workflows cover variant calling, read alignment, taxonomic
profiling, and metagenomic classification, producing different relationships
between staged input size, measured consumed data, and peak memory.

% Bowtie2 is a separate custom live workload with 6,000 logical alignment tasks
% per run. It is excluded from the offline training/test split so that the live
% case study evaluates a cold run, persistence of completed observations, model
% retraining, and a fresh warm rerun on a workflow outside the offline cohort.

\subsection{\toolname Model Instantiation}

The memory-allocation experiment evaluates two conditions on the same tasks.
\emph{Access} is the peak-history baseline: it uses the
launch-time static vector and causal peak-memory history. \emph{\toolname} is the
consumption-aware allocator: it additionally uses causal consumption history
from earlier audited tasks and a prelaunch estimate of the current task's
consumption. Their legal input vectors are
\begin{align}
\mathbf{x}_t^{\mathrm{Access}}
   &= [\mathbf{s}_t,\mathbf{h}_{m,t}],\label{eq:access-view}\\
\hat c_t &= g(\mathbf{s}_t,\mathbf{h}_{m,t}),\nonumber\\
\mathbf{x}_t^{\mathrm{\toolname{}}}
   &= [\mathbf{s}_t,\mathbf{h}_{m,t},\mathbf{h}_{c,t},\hat c_t].
\label{eq:camp-view}
\end{align}
Thus, neither condition can observe the current task's measured $c_t$ or $y_t$
before launch. Table~\ref{tab:model-instantiation} maps these inputs to the
outputs produced by each condition using the notation from Section~III.

\begin{table}[t]
\caption{Access and \toolname model instantiation.}
\label{tab:model-instantiation}
\centering
\scriptsize
\setlength{\tabcolsep}{2.5pt}
\renewcommand{\arraystretch}{1.08}
\begin{tabular}{@{}p{0.44\columnwidth}p{0.47\columnwidth}@{}}
\toprule
Input parameters & Model outputs \\
\midrule
Access: $\mathbf{x}_t^{\mathrm{Access}}$ &
$\{Q_q^{\mathrm{Access}}(t)\}_{q\in\mathcal Q}$ and calibrated request
$\hat y_t^{\mathrm{Access}}$ \\
\toolname: $\mathbf{x}_t^{\mathrm{\toolname{}}}$ &
$\{Q_q^{\mathrm{\toolname{}}}(t)\}_{q\in\mathcal Q}$ and calibrated request
$\hat y_t^{\mathrm{\toolname{}}}$ \\
\bottomrule
\end{tabular}
\end{table}

For each condition, the global point model and the matching workflow-process
upper-tail model produce the quantile estimates listed in
Table~\ref{tab:model-instantiation}; an unseen workflow-process pair uses the
global upper-tail fallback. The calibration procedure in
Section~\ref{sec:calibration} evaluates the prediction and peak-history
candidates and converts the selected base into the submitted memory request.
Both conditions use the same task rows, peak-RSS target, quantile model family,
and leakage-control protocol; each is fitted and calibrated using development
data only. The next subsection describes this common model-training and
calibration procedure.

  \subsection{Training and Calibration Protocol}

The 33,516-task cohort is divided chronologically, with the earlier 28,490
tasks used for development and the latest 5,026 tasks reserved for held-out
evaluation. Related task-instance, staged-input, and sample groups remain in
the same partition. Access and \toolname use the same development and test tasks,
peak-RSS target, model family, and preprocessing. Their legal inputs are defined
above, and each condition is fitted and calibrated independently. Peak RSS is
modelled as $\log(1+y_t)$, with $y_t$ expressed in MiB, while the consumption
estimator uses $\log(1+c_t)$, with $c_t$ expressed in bytes.

Five-fold grouped out-of-fold prediction is performed on the development
tasks. Consequently, each development prediction is produced by models that
were not trained on that task or its associated split group. These out-of-fold
predictions support calibration and policy selection without exposing a row to
its own target during prediction. Automated checks enforce identifier
uniqueness, deterministic order, feature formulas, system compatibility, and
exclusion of current post-run fields.

A deterministic task-identifier hash separates the out-of-fold development
predictions into 17,210 residual-reference rows and 11,280 policy-selection
rows. The residual-reference rows construct the condition-specific residual
distributions used for calibration. Candidate allocation policies are then
evaluated on the policy-selection rows. Among candidates satisfying at
least 99\% overall coverage and 97.5\% coverage for every workflow, each
condition selects the policy with the smallest request-to-peak ratio.

After policy selection, the models are refitted on all 28,490 development
tasks, and their parameters and selected policies are frozen before held-out
evaluation. The 5,026 test tasks are processed sequentially. A task's own
outcome is unavailable when its request is produced; once revealed, it may
update causal history and residual state for later tasks. No test outcome is
used to refit model weights or reselect the allocation policy. The offline
replay approximates recorded completion order, whereas live \toolname{} admits
only outcomes completed before decision time $d_t$.

Together, these controls support four research questions" %that progress from
% measurement, through memory allocation and live deployment, to external
% comparison and selective feedback:

\begin{enumerate}[leftmargin=*,label=\textbf{RQ\arabic*:},itemsep=1pt,topsep=2pt]
\item \textbf{eBPF and mmap signal.} Compared with STRACE, how much
explicit-read and memory-mapped page-fault activity does eBPF capture, and
what runtime overhead does online full auditing introduce?
\item \textbf{Consumption-aware allocation.} Does \toolname reduce
counterfactual first-attempt underallocation and runtime-weighted unused
memory relative to Access, and what point-error and reservation-cost
tradeoffs result?
% \item \textbf{Live Bowtie2 feedback.} In fresh cold and warm executions, can
% CAMP allocate memory before each task, learn from completed observations, and
% reduce warm-run memory cost while maintaining first-attempt coverage?
\item \textbf{Comparison with Sizey.} On identical offline task assignments,
how does \toolname compare with unchanged Sizey in underallocation, point error,
and reserved-memory cost?
\item \textbf{Selective auditing.} At equal eBPF audit counts, do \toolname{}'s
gates identify more underallocated, memory-heavy, and
consumption-informative tasks than random selection?
\item \textbf{Live Bowtie2 feedback.} In fresh cold and warm executions, can
\toolname{} allocate memory before each task, learn only from completed
observations, and reduce warm-run reservation cost while maintaining high
first-attempt coverage?
\end{enumerate}

\toolname{} trains a global consumption ensemble and, for each condition, a
global point ensemble, workflow-process upper-tail ensembles, and a global
upper-tail fallback. Each peak-memory quantile uses three fixed-seed LightGBM
models with the Access and \toolname{} inputs in
Equations~\ref{eq:access-view} and~\ref{eq:camp-view}. Scope-pool construction
dominates
cost: key-matched
scopes are $O(N)$, while $I_2$'s nearest-neighbor search requires indexing at
scale. At inference, \toolname{} evaluates the consumption estimator, global
point model, and matching workflow-process upper-tail model, using the global
upper-tail fallback for an unseen pair.

\subsection{Evaluation Protocol and Metrics}

For first request $\hat y_t$, recorded peak RSS $y_t$, q50 prediction
$\tilde y_t$, and runtime $\tau_t$, the common allocation metrics are
\begin{align*}
u_t&=\mathbf1[\hat y_t<y_t], &
\mathrm{Cov}&=1-\frac{1}{N}\sum_tu_t,\\
\kappa&=\frac{\sum_t \hat y_t}{\sum_t y_t}, &
W&=\sum_t\tau_t(\hat y_t-y_t)_+,\\
\mathrm{MdAPE}&=100\,\underset{t}{\mathrm{median}}
 \frac{|\tilde y_t-y_t|}{y_t}.&&
\end{align*}
Here $u_t$ marks a counterfactual underallocation, $\kappa$ is the aggregate
request-to-peak ratio, and $W$ is unused memory-time, reported in GiB-hours.
We also report total requested memory.

\noindent\textbf{RQ1.} Matched no-audit, full-STRACE, and full-eBPF executions
report explicit-read and mmap-fault signal shares, absolute mmap-fault bytes,
and runtime overhead $100(T_w^a/T_w^0-1)$ relative to no audit.

\noindent\textbf{RQ2.} Access and \toolname{} are compared using held-out
underallocation counts across chronological grouped 30/70, 50/50, 70/30,
85/15, and 90/10 splits. Seeds 1996--2000 refit all models and policies at each
split. For \toolname{}, underallocations per 1,000 tasks, the five-seed mean,
$t$-based 95\% confidence interval, and coefficient of variation measure seed
sensitivity on identical holdout rows. For the primary 85/15 split, total
requested memory and $W$ are also reported.

\noindent\textbf{RQ3.} Sizey and \toolname{} receive identical tasks under
Sizey's 9,141/21,337-task per-process split. Sizey's process-size eligibility
rule removes approximately 3,000 records. Underallocation count, coverage, q50
MdAPE, $\kappa$, and $W$ are reported, with $\kappa$ and $W$ normalized to
Sizey.

\noindent\textbf{RQ4.} Similarly 9,141-task development set selects and freezes the
three-gate policy before replaying 21,337 later tasks. At budget $B$, recall of
target $z_t$ is
\begin{equation}
\operatorname{Recall}_z(B)=
\frac{\sum_t q_t(B)z_t}{\sum_t z_t}.
\label{eq:audit-recall}
\end{equation}
Risk-target and high-information-target recall are compared with uniform and
process-stratified random selection at equal counts. Budgets of 5\%, 10\%,
20\%, 30\%, and 50\% use 1,000 repetitions; random-baseline shading gives
empirical 95\% intervals.

\noindent\textbf{RQ5.} Two fresh 6,000-task Bowtie2 executions use the same
workflow structure, inputs, worker pool, \toolname{} allocation policy, and
retry policy. The cold run starts without Bowtie2 history; the warm run starts
from the terminal model and completed-history state of the cold run. The runs
use separate Nextflow work directories, full eBPF auditing, and no
\texttt{-resume}. We report aggregate observed peak and first-requested memory,
first-attempt coverage, unused memory-time, allocation-margin buckets,
successful tasks, and recorded first-attempt OOMs.
